% META: {"id": "TCOMPL-MODELES-EVOLUTION-QCM", "chapitre": "TCOMPL-MODELES-EVOLUTION", "type_objet": "qcm", "genere_depuis": "chapitres/TCOMPL-MODELES-EVOLUTION/qcm/TCOMPL-MODELES-EVOLUTION-QCM.json", "status": "generated"}
% Fichier genere par scripts/build_qcm_tex.py — ne pas editer a la main.

\section*{\textcolor{chapcolor}{\MakeUppercase{Modèles d'evolution discrète et continue}}}

\begin{center}
\textit{Pour chaque question, une seule réponse est exacte.}
\end{center}

\begin{enumerate}

\item[] \textbf{Capacité C1}

\item \textbf{[Q1]} Une suite arithmétique de raison r et premier terme $u_0$ a pour terme general :
  \begin{enumerate}[label=\Alph*.]
    \item $u_n = u_0 + n r$
    \item $u_n = u_0 \times r^n$
    \item $u_n = u_0 + r^n$
    \item $u_n = u_0 n r$
  \end{enumerate}

\item[] \textbf{Capacité C2}

\item \textbf{[Q2]} Pour $0<q<1$, la somme $1+q+q^2+\cdots$ vaut :
  \begin{enumerate}[label=\Alph*.]
    \item $1/(1-q)$
    \item $1-q$
    \item $1/q$
    \item $+\infty$
  \end{enumerate}

\item[] \textbf{Capacité C5}

\item \textbf{[Q3]} L'équation différentielle $y' = a y$ a pour solutions les fonctions :
  \begin{enumerate}[label=\Alph*.]
    \item $f(x) = a x + C$
    \item $f(x) = C e^{a x}$
    \item $f(x) = C e^{-a x}$
    \item $f(x) = C x^a$
  \end{enumerate}

\item \textbf{[Q4]} Le modèle de Malthus suppose un taux de croissance :
  \begin{enumerate}[label=\Alph*.]
    \item Decroissant
    \item Nul
    \item Constant et proportionnel a la population
    \item Logistique
  \end{enumerate}

\item[] \textbf{Capacité C4}

\item \textbf{[Q5]} Une solution constante de $u_{n+1}=0{,}8u_n+10$ vaut :
  \begin{enumerate}[label=\Alph*.]
    \item $10$
    \item $40$
    \item $80$
    \item $50$
  \end{enumerate}

\end{enumerate}

% NEXUS-QCM-TEACHER-ONLY-BEGIN
\ifnxVersionProfesseur
\clearpage
\section*{Clé de correction — réservée au professeur}

\begin{center}
\begin{tabular}{lll}
\hline
Question & Capacité & Réponse exacte \\
\hline
Q1 & C1 & \textbf{A} \\
Q2 & C2 & \textbf{A} \\
Q3 & C5 & \textbf{B} \\
Q4 & C5 & \textbf{C} \\
Q5 & C4 & \textbf{D} \\
\hline
\end{tabular}
\end{center}

\subsection*{Diagnostic des réponses erronées}

\begin{itemize}
  \item \textbf{Q1}
  \begin{itemize}
    \item \textbf{B} — Formule d'une suite géométrique. Une suite arithmétique de raison r vérifie u\_n=u\_0+nr. \emph{Renvoi : C1.}
    \item \textbf{C} — Formule incorrecte. Une suite arithmétique de raison r vérifie u\_n=u\_0+nr. \emph{Renvoi : C1.}
    \item \textbf{D} — Produit au lieu d'addition. Une suite arithmétique de raison r vérifie u\_n=u\_0+nr. \emph{Renvoi : C1.}
  \end{itemize}
  \item \textbf{Q2}
  \begin{itemize}
    \item \textbf{B} — La différence $1-q$ apparaît au dénominateur après multiplication et soustraction de la série, pas comme résultat. Pour 0<q<1, la somme géométrique infinie 1+q+q\textasciicircum{}2+... vaut 1/(1-q). \emph{Renvoi : C2.}
    \item \textbf{C} — La valeur $1/q$ ne vérifie pas l'identité S=1+qS obtenue à partir de la série. Pour 0<q<1, la somme géométrique infinie 1+q+q\textasciicircum{}2+... vaut 1/(1-q). \emph{Renvoi : C2.}
    \item \textbf{D} — Comme la raison est strictement comprise entre 0 et 1, les sommes partielles convergent vers une limite finie. Pour 0<q<1, la somme géométrique infinie 1+q+q\textasciicircum{}2+... vaut 1/(1-q). \emph{Renvoi : C2.}
  \end{itemize}
  \item \textbf{Q3}
  \begin{itemize}
    \item \textbf{A} — Solution de y' = a. Toutes les solutions de y'=ay sont y(x)=C exp(ax), avec C réel. \emph{Renvoi : C5.}
    \item \textbf{C} — Erreur de signe dans l'exposant. Toutes les solutions de y'=ay sont y(x)=C exp(ax), avec C réel. \emph{Renvoi : C5.}
    \item \textbf{D} — Fonction puissance au lieu d'exponentielle. Toutes les solutions de y'=ay sont y(x)=C exp(ax), avec C réel. \emph{Renvoi : C5.}
  \end{itemize}
  \item \textbf{Q4}
  \begin{itemize}
    \item \textbf{A} — Hypothèse du modèle de Verhulst. Le modèle de Malthus s'écrit y'=ay : le coefficient de croissance a est constant et la variation est proportionnelle à la population. \emph{Renvoi : C5.}
    \item \textbf{B} — Population constante. Le modèle de Malthus s'écrit y'=ay : le coefficient de croissance a est constant et la variation est proportionnelle à la population. \emph{Renvoi : C5.}
    \item \textbf{D} — Modèle avec capacité d'accueil. Le modèle de Malthus s'écrit y'=ay : le coefficient de croissance a est constant et la variation est proportionnelle à la population. \emph{Renvoi : C5.}
  \end{itemize}
  \item \textbf{Q5}
  \begin{itemize}
    \item \textbf{A} — La constante additive 10 n'est pas le point fixe, car $0{,}8*10+10=18$. Une solution constante de u\_(n+1)=0,8u\_n+10 vérifie ell=0,8ell+10, donc ell=50. \emph{Renvoi : C4.}
    \item \textbf{B} — La valeur 40 ne satisfait pas l'équation fixe : $0{,}8*40+10=42$. Une solution constante de u\_(n+1)=0,8u\_n+10 vérifie ell=0,8ell+10, donc ell=50. \emph{Renvoi : C4.}
    \item \textbf{C} — La valeur 80 ne satisfait pas l'équation fixe : $0{,}8*80+10=74$. Une solution constante de u\_(n+1)=0,8u\_n+10 vérifie ell=0,8ell+10, donc ell=50. \emph{Renvoi : C4.}
  \end{itemize}
\end{itemize}
\fi
% NEXUS-QCM-TEACHER-ONLY-END
